\section{Introduction}
\label{sec:introduction}
The human gait cycle—the sequence of movements occurring between two successive foot contacts of the same limb—offers critical insights into motor coordination, 
balance, and neurological function. In clinical practice, accurate assessment of the gait cycle is essential for identifying pathological asymmetries, especially 
in post-stroke patients, where hemiparesis often leads to disrupted stance and swing phase dynamics \cite{ref1, ref2}.

Traditional gait analysis systems using optical motion capture and force plates provide detailed temporal and kinematic data but remain impractical for routine 
use due to high costs, space requirements, and operational complexity \cite{ref3, ref4}. These limitations hinder frequent assessments in stroke rehabilitation, 
where ongoing monitoring of gait asymmetry is vital to tailor therapy and track recovery progress \cite{ref5}.

Recent advancements in wearable sensors—particularly inertial measurement units (IMUs)—have enabled portable, cost-effective alternatives for gait cycle analysis 
outside laboratory settings \cite{ref6, ref7}. When paired with machine learning, IMU data can be used to extract interpretable features that reflect clinically 
relevant gait events.

This study introduces a lightweight and interpretable pipeline that classifies gait cycle asymmetries using bilateral gyroscopic data. By focusing on extrema-based 
statistical features and employing subject-wise validation, the proposed approach balances accuracy and computational efficiency—supporting real-world deployment 
in embedded or web-based rehabilitation tools.

\textbf{The main contributions are:}
\begin{itemize}
    \item A minimal-feature extraction method using bilateral gyroscope signals for gait cycle asymmetry detection.
    \item A comparative evaluation of traditional machine learning models using leave-one-out cross-validation (LOOCV) across subjects.
    \item Demonstration of high classification performance (AUC $>$ 0.95) with reduced computational complexity.
\end{itemize}

